\documentclass[11pt]{article}
\usepackage[utf8]{inputenc}
\usepackage[T1]{fontenc}
\usepackage{amsmath,amssymb,amsfonts}
\usepackage{booktabs}
\usepackage{graphicx}
\usepackage{hyperref}
\usepackage{xcolor}
\usepackage{multirow}
\usepackage{tabularx}
\usepackage[margin=1in]{geometry}
\usepackage{natbib}
\usepackage{array}

\newcommand{\lzstar}{$l_{z}^{*}$}
\newcommand{\lzm}{$l_{z}^{m}$}

\title{On the Fundamental Limits of Binary-Response Contamination Detection:\\Item Response Theory Person-Fit Statistics Applied to Large Language Models}

\author{Anonymous}
\date{}

\begin{document}
\maketitle

% ============================================================
% ABSTRACT
% ============================================================
\begin{abstract}
Benchmark contamination threatens the reliability of Large Language Model (LLM) evaluation, yet all existing detection methods require privileged access beyond binary response patterns---training data, token probabilities, model weights, or live inference.
We investigate whether Item Response Theory (IRT) person-fit statistics, developed for detecting exam cheating in psychometrics, can operate at this minimal-information operating point.
Through simulation, controlled contamination experiments, and a leaderboard-scale audit, we establish three findings:
(1)~LLM response patterns \emph{universally} violate IRT assumptions---local independence and unidimensionality---producing extreme person-fit scores ($l_z^* \gg 0$) even for uncontaminated models, with 93\% of model--benchmark cells flagged under absolute thresholds;
(2)~Supervised fine-tuning (SFT) contamination induces ability shifts inseparable from genuine improvement in binary responses, yielding inconsistent $\Delta l_z^*$ direction across contamination levels (Cohen's $d = 0.25$, AUROC $= 0.585$ at best) and reducing to chance after reinforcement learning (AUROC $= 0.53$);
(3)~a theoretical power analysis reveals an inverted-U detection gap---power reaches its nadir ($\leq 0.12$) precisely at moderate contamination levels (25--40\%) where detection matters most.
These results establish that binary-response patterns contain insufficient information for reliable contamination detection, quantify the detection power gained at each information-access level (binary $\rightarrow$ re-test: $+0.235$ AUROC; binary $\rightarrow$ token probabilities: $+0.195$), and provide information-theoretic justification for gray-box detection methods.
\end{abstract}

% ============================================================
% 1. INTRODUCTION
% ============================================================
\section{Introduction}

Benchmark contamination---the presence of evaluation data in model training corpora---is a growing threat to the validity of LLM evaluation.
As training datasets scale to trillions of tokens and benchmark items circulate widely online, the probability of inadvertent or intentional data leakage increases.
Contaminated benchmark scores inflate perceived model capability, distort leaderboard rankings, and mislead downstream deployment decisions.

Existing detection methods span a wide spectrum of information requirements.
N-gram overlap methods require access to training data \citep{brown2020gpt3}.
Membership Inference Attacks (MIA) and Min-K\%++ require token-level log probabilities \citep{shi2024minkpp}.
ConStat requires live model access for rephrased re-testing \citep{dekoninck2024constat}.
Internal probing methods require model weights.
Critically, \citet{wang2025mia} demonstrate that all MIA-based methods fail under post-reinforcement-learning contamination, highlighting the fragility of even token-level approaches.

A natural question arises: \emph{what can be detected from binary response patterns alone?}
This is the minimal-information operating point on the contamination detection frontier---requiring only the correct/incorrect outcomes that benchmarks already report.
If effective, such a method could audit any model with published benchmark scores, including the 7,200+ models on the Open LLM Leaderboard, without requiring model access.

Item Response Theory (IRT) person-fit statistics, developed over decades for detecting exam cheating in educational testing \citep{meijer2001person,snijders2001}, provide a principled framework for this operating point.
The $l_z^*$ statistic \citep{snijders2001} measures how well a respondent's pattern of correct and incorrect answers matches the IRT model's predictions given their estimated ability.
A contaminated model---like an examinee with item preknowledge---should exhibit aberrant response patterns detectable via person-fit analysis.

In this paper, we conduct the first systematic investigation of IRT person-fit statistics for LLM contamination detection and establish that this approach faces \emph{fundamental} limitations.
Our investigation proceeds through three complementary analyses: theoretical power characterization, controlled contamination experiments with ground-truth labels, and a leaderboard-scale empirical audit.
The results converge on a single conclusion: binary response patterns contain insufficient information to reliably separate contamination from genuine ability, particularly after SFT-induced knowledge transfer.

\paragraph{Contributions.}
\begin{enumerate}
    \item \textbf{Universal misfit characterization.} We demonstrate that LLM responses systematically violate IRT model assumptions (local independence, unidimensionality), producing extreme $l_z^*$ values that render absolute detection thresholds useless. Across 42 leaderboard models, 93\% of model--benchmark cells are flagged, with a cross-benchmark Spearman correlation of only $\rho = 0.18$.

    \item \textbf{SFT contamination--ability confound.} Through controlled experiments on two models (Qwen3.5-9B, Gemma-4-31B) across six contamination levels, we show that SFT contamination produces $\Delta l_z^*$ effects that are small (Cohen's $d = 0.25$), direction-inconsistent (60\% cross-model agreement), and model-capacity-dependent. The best AUROC achieved is 0.585---functionally near random.

    \item \textbf{Detection information hierarchy.} By comparing IRT-FIT against ConStat (AUROC = 0.82) and Min-K\%++ (AUROC = 0.78), we quantify the detection power gained at each information-access level, establishing that moving from binary responses to rephrased re-testing yields $+0.235$ AUROC---a concrete justification for gray-box methods.

    \item \textbf{Theoretical power lower bound.} We derive a closed-form power function from the Snijders variance formula, revealing an inverted-U detection gap where power reaches its minimum at moderate contamination ($\sim$35\%), precisely the regime where detection is most needed.
\end{enumerate}

% ============================================================
% 2. BACKGROUND
% ============================================================
\section{Background}

\subsection{Item Response Theory}

IRT models the probability of a correct response as a function of the respondent's latent ability $\theta$ and item characteristics.
The four-parameter logistic (4PL) model specifies:
\begin{equation}
    P(X_{ij} = 1 \mid \theta_j) = c_i + (d_i - c_i) \cdot \frac{1}{1 + e^{-a_i(\theta_j - b_i)}}
\end{equation}
where $a_i$ is the discrimination parameter, $b_i$ is the difficulty, $c_i$ is the lower asymptote (guessing), and $d_i$ is the upper asymptote (feasibility).
We adopt 4PL item parameters from the PSN-IRT benchmark calibration, which provides pre-estimated parameters for major LLM benchmarks.

\subsection{Person-Fit Statistics}

Person-fit statistics evaluate whether an individual's response pattern is consistent with the IRT model given their estimated ability $\hat{\theta}$.
The $l_z^*$ statistic \citep{snijders2001} is the standardized log-likelihood:
\begin{equation}
    l_z^* = \frac{l_0 - E(l_0)}{\sqrt{\text{Var}(l_0)}}
\end{equation}
where $l_0 = \sum_{i=1}^{n} [x_i \log P_i + (1-x_i) \log(1 - P_i)]$ is the log-likelihood of the response vector.
Under the null hypothesis of model fit, $l_z^* \sim \mathcal{N}(0, 1)$ asymptotically.

Negative $l_z^*$ values indicate aberrant patterns (unexpected errors on easy items or unexpected successes on hard items), while positive values indicate over-consistent patterns (responses more predictable than the model expects).
In the classical testing literature, examinees with item preknowledge typically produce negative $l_z^*$ (they answer leaked items correctly regardless of difficulty, creating a mismatch with their estimated ability on non-leaked items).

\subsection{Contamination Detection Methods}

We position IRT-FIT within a detection information hierarchy:
\begin{itemize}
    \item \textbf{Level 0 (binary responses):} IRT-FIT---requires only correct/incorrect patterns.
    \item \textbf{Level 1 (re-test access):} ConStat \citep{dekoninck2024constat}---requires rephrased benchmark items and live model access.
    \item \textbf{Level 2 (token probabilities):} Min-K\%++ \citep{shi2024minkpp}---requires token-level log probabilities.
    \item \textbf{Level 3+ (logits/weights):} MIA variants, internal probing---require deep model access.
\end{itemize}

% ============================================================
% 3. METHOD: IRT-FIT FRAMEWORK
% ============================================================
\section{IRT-FIT Detection Framework}

\subsection{Framework Overview}

Given a model's binary response vector $\mathbf{x} = (x_1, \ldots, x_n)$ on a benchmark with pre-calibrated 4PL item parameters, IRT-FIT proceeds as follows:
(1)~estimate the model's ability $\hat{\theta}$ via maximum likelihood;
(2)~compute the person-fit statistic $l_z^*$;
(3)~compare against a detection threshold.

\subsection{Absolute vs.\ Relative Detection}

An absolute threshold (e.g., $|l_z^*| > 1.96$) assumes the IRT model fits LLM response data.
As we demonstrate in Section~\ref{sec:leaderboard}, this assumption fails universally for LLMs.
We therefore introduce a \emph{relative} detection approach: given a clean baseline model's ${l_z^*}_{\text{clean}}$, we define the contamination signal as $\Delta l_z^* = {l_z^*}_{\text{contam}} - {l_z^*}_{\text{clean}}$ and test whether this shift exceeds natural variation.

\subsection{Multiscale Aggregation}

Following \citet{hong2021multiscale}, we compute a multiscale statistic $l_z^m$ that aggregates standardized person-fit scores across benchmarks, aiming to increase power through cross-benchmark evidence combination.

% ============================================================
% 4. THEORETICAL ANALYSIS
% ============================================================
\section{Theoretical Power Analysis}
\label{sec:theory}

\subsection{Power Function Derivation}

We derive the detection power as a function of the number of items $n$, contamination proportion $\pi$, and the ability--difficulty gap.
Starting from the Snijders variance formula $\text{Var}(l_z^*) \approx 1 + O(1/n)$ under the null, we model the non-null distribution under contamination by assuming contaminated items are answered correctly with probability 1, independent of item difficulty.

The resulting power function exhibits an \emph{inverted-U} shape as a function of contamination level $\pi$:
at low $\pi$, the contamination signal is too weak relative to sampling noise;
at high $\pi$, the response pattern becomes over-consistent (all correct), which under re-estimated $\hat{\theta}$ appears as a legitimately high-ability examinee, collapsing the aberrance signal.
The theoretical nadir occurs at $\pi \approx 0.33$, closely matching the empirical nadir at $\pi = 0.35$ (Table~\ref{tab:inverted-u}).

\subsection{Theoretical vs.\ Simulated Power}

Table~\ref{tab:theory-vs-sim} compares theoretical and simulated power at 15\% contamination.
The theoretical curve systematically overestimates power by 15--30\%, because the Snijders formula assumes local independence among item responses---an assumption violated by the correlated contamination pattern.
The average $R^2$ between theoretical and simulated power curves is 0.61, sufficient to capture the qualitative inverted-U shape but too loose for practical calibration (the closed-form bound overestimates by a factor of $\sim$2$\times$ at moderate contamination).

\begin{table}[t]
\centering
\caption{Theoretical vs.\ simulated detection power at 15\% contamination under the 4PL model.}
\label{tab:theory-vs-sim}
\begin{tabular}{lrrrc}
\toprule
Benchmark & $n_\text{items}$ & Power$_\text{theory}$ & Power$_\text{sim}$ & Ratio \\
\midrule
MMLU       & 14,042 & 0.996 & 0.979 & 1.02 \\
HellaSwag  & 10,042 & 0.880 & 0.733 & 1.20 \\
GSM8K      &  1,319 & 0.720 & 0.578 & 1.25 \\
ARC-C      &  1,172 & 0.180 & 0.129 & 1.40 \\
HumanEval  &    164 & 0.360 & 0.243 & 1.48 \\
\bottomrule
\end{tabular}
\end{table}

\subsection{The Inverted-U Detection Gap}

\begin{table}[t]
\centering
\caption{Inverted-U characterization across benchmarks. The nadir contamination level denotes where detection power is minimized.}
\label{tab:inverted-u}
\begin{tabular}{lrcr}
\toprule
Benchmark & Nadir (\%) & Nadir Power & $R^2_\text{theory vs.\ sim}$ \\
\midrule
MMLU      & 35 & 0.12 & 0.68 \\
HellaSwag & 33 & 0.08 & 0.62 \\
GSM8K     & 30 & 0.05 & 0.55 \\
ARC-C     & 32 & 0.03 & 0.51 \\
HumanEval & 28 & 0.04 & 0.53 \\
\bottomrule
\end{tabular}
\end{table}

The inverted-U reveals a structural weakness of binary-response detection: power reaches its minimum at 25--40\% contamination (Table~\ref{tab:inverted-u}), precisely the moderate contamination regime most relevant to real-world leaderboard gaming.
At the nadir, even MMLU (14,042 items) achieves only 0.12 power.
Smaller benchmarks fare worse: ARC-C (1,172 items) has nadir power of 0.03---essentially undetectable.

% ============================================================
% 5. SIMULATION STUDY
% ============================================================
\section{Simulation Study}
\label{sec:simulation}

We conduct a comprehensive simulation using calibrated 4PL item parameters to establish the theoretical detection ceiling under idealized conditions (perfect IRT model fit).

\subsection{Setup}

For each benchmark, we simulate response vectors under the 4PL model at contamination levels $\pi \in \{0.15, 0.25, 0.35, 0.50, 0.60, 0.75, 0.90, 1.00\}$.
Contaminated items are set to $x_i = 1$ (correct) regardless of ability.
We compute two-sided detection rates: negative-side ($l_z^* < -1.96$, aberrant patterns) and positive-side ($l_z^* > 1.96$, over-consistent patterns).

\subsection{Results}

\begin{table}[t]
\centering
\caption{Two-sided detection power across contamination levels (simulation). Values denote the proportion of simulated examinees detected at $\alpha = 0.05$. At low contamination, detection is primarily negative-side (aberrant); at high contamination, primarily positive-side (over-consistent). The crossover for MMLU occurs at $\pi \approx 0.28$.}
\label{tab:two-sided-power}
\resizebox{\textwidth}{!}{%
\begin{tabular}{lrcccccccc}
\toprule
Benchmark & $n$ & 15\%$_\text{neg}$ & 25\%$_\text{neg}$ & 35\%$_\text{comb}$ & 50\%$_\text{pos}$ & 60\%$_\text{pos}$ & 75\%$_\text{pos}$ & 90\%$_\text{pos}$ & 100\%$_\text{pos}$ \\
\midrule
MMLU      & 14,042 & 0.979 & 0.72 & 0.12 & 0.31 & 0.42 & 0.61 & 0.78 & 0.85 \\
HellaSwag & 10,042 & 0.733 & 0.48 & 0.08 & 0.18 & 0.24 & 0.28 & 0.52 & 0.64 \\
GSM8K     &  1,319 & 0.578 & 0.31 & 0.05 & 0.09 & 0.12 & 0.15 & 0.22 & 0.31 \\
ARC-C     &  1,172 & 0.129 & 0.06 & 0.03 & 0.04 & 0.05 & 0.06 & 0.09 & 0.14 \\
HumanEval &    164 & 0.243 & 0.11 & 0.04 & 0.05 & 0.06 & 0.07 & 0.10 & 0.13 \\
\bottomrule
\end{tabular}%
}
\end{table}

Table~\ref{tab:two-sided-power} presents the full power landscape.
At 15\% contamination, MMLU achieves 0.979 power via negative-side detection---contaminated items create aberrant ``easy items answered correctly despite low estimated ability'' patterns.
However, as contamination increases to 35\%, negative-side power decays (the ability estimate $\hat{\theta}$ rises, rationalizing more correct answers) while positive-side power has not yet accumulated, creating the inverted-U nadir at power $= 0.12$.

Two key observations emerge.
First, power is strongly dependent on $n_\text{items}$: MMLU (14K items) achieves useful power at extreme contamination levels, while HumanEval (164 items) never exceeds 0.25 even at 100\% contamination.
Second, the simulation assumes perfect IRT model fit---real LLM data will exhibit additional model misfit that further degrades detection power.

% ============================================================
% 6. CONTROLLED CONTAMINATION EXPERIMENTS
% ============================================================
\section{Controlled Contamination Experiments}
\label{sec:controlled}

Simulation establishes theoretical ceilings. To evaluate detection in realistic conditions, we conduct controlled contamination experiments with ground-truth labels.

\subsection{Setup}

We fine-tune two models---Qwen3.5-9B and Gemma-4-31B---using QLoRA SFT on MMLU items at contamination levels $\pi \in \{0, 0.05, 0.10, 0.15, 0.25, 0.50\}$.
Each condition is repeated three times with different random seeds.
We measure $\Delta l_z^* = {l_z^*}_{\text{contam}} - {l_z^*}_{\text{clean}}$ and compute AUROC for discriminating contaminated from clean models.

\subsection{Pilot Results}

An initial pilot on Qwen3.5-9B $\times$ MMLU (Table~\ref{tab:pilot}) revealed the core obstacle: SFT contamination simultaneously changes the model's ability $\theta$ and its response pattern, and the 4PL model's re-estimation of $\hat{\theta}$ absorbs much of the contamination signal.

\begin{table}[t]
\centering
\caption{Pilot results: Qwen3.5-9B $\times$ MMLU. The $\Delta l_z^*$ direction is inconsistent across contamination levels, and neither shift is statistically significant.}
\label{tab:pilot}
\begin{tabular}{lccccc}
\toprule
Condition & Accuracy & $\theta$ & $l_z^*$ & $\Delta l_z^*$ & $p$-value \\
\midrule
Clean (0\%)  & 0.71 & 1.05 & 18.4 & ---   & --- \\
15\% contam. & 0.76 & 1.28 & 21.1 & $+$2.7  & 0.18 \\
50\% contam. & 0.84 & 1.72 & 16.9 & $-$1.5  & 0.42 \\
\bottomrule
\end{tabular}
\end{table}

At 15\% contamination, $\Delta l_z^* = +2.7$ ($p = 0.18$): the contamination induces a slight increase in person-fit misfit, but the effect is not significant.
At 50\%, $\Delta l_z^* = -1.5$ ($p = 0.42$): the direction \emph{reverses}---massive contamination produces an over-consistent pattern that the model interprets as legitimately high ability ($\hat{\theta} = 1.72$).
The direction inconsistency and large run-to-run variance (SD $\approx$ 7.5) led to a NO-GO decision for the detection framing, motivating the pivot to impossibility characterization.

\subsection{Full Controlled Experiment}

\begin{table}[t]
\centering
\caption{$\Delta l_z^*$ across contamination levels for Qwen3.5-9B on MMLU. No contamination level produces a statistically significant shift.}
\label{tab:ctrl-qwen}
\begin{tabular}{lccccccr}
\toprule
Contam.\,\% & Acc. & $\Delta$Acc. & $\theta$ & $\Delta\theta$ & $l_z^*$ & $\Delta l_z^*$ & $p$ \\
\midrule
0\% (clean) & 0.71 & ---    & 1.05 & ---    & 18.4 & ---    & --- \\
5\%         & 0.73 & $+$0.02 & 1.12 & $+$0.07 & 19.2 & $+$0.8  & 0.64 \\
10\%        & 0.75 & $+$0.04 & 1.21 & $+$0.16 & 20.3 & $+$1.9  & 0.31 \\
15\%        & 0.76 & $+$0.05 & 1.28 & $+$0.23 & 21.1 & $+$2.7  & 0.18 \\
25\%        & 0.80 & $+$0.09 & 1.49 & $+$0.44 & 19.6 & $+$1.2  & 0.55 \\
50\%        & 0.84 & $+$0.13 & 1.72 & $+$0.67 & 16.9 & $-$1.5  & 0.42 \\
\bottomrule
\end{tabular}
\end{table}

\begin{table}[t]
\centering
\caption{$\Delta l_z^*$ across contamination levels for Gemma-4-31B on MMLU. The larger model absorbs contamination more completely, with $\Delta l_z^*$ trending negative at higher contamination.}
\label{tab:ctrl-gemma}
\begin{tabular}{lccccccr}
\toprule
Contam.\,\% & Acc. & $\Delta$Acc. & $\theta$ & $\Delta\theta$ & $l_z^*$ & $\Delta l_z^*$ & $p$ \\
\midrule
0\% (clean) & 0.78 & ---    & 1.41 & ---    & 22.7 & ---    & --- \\
5\%         & 0.80 & $+$0.02 & 1.50 & $+$0.09 & 22.4 & $-$0.3  & 0.88 \\
10\%        & 0.82 & $+$0.04 & 1.59 & $+$0.18 & 23.3 & $+$0.6  & 0.78 \\
15\%        & 0.83 & $+$0.05 & 1.66 & $+$0.25 & 23.8 & $+$1.1  & 0.62 \\
25\%        & 0.86 & $+$0.08 & 1.82 & $+$0.41 & 21.9 & $-$0.8  & 0.71 \\
50\%        & 0.91 & $+$0.13 & 2.15 & $+$0.74 & 19.5 & $-$3.2  & 0.22 \\
\bottomrule
\end{tabular}
\end{table}

Tables~\ref{tab:ctrl-qwen} and~\ref{tab:ctrl-gemma} present the full results.
Three patterns emerge:

\paragraph{Direction inconsistency.}
$\Delta l_z^*$ is positive at low-to-moderate contamination (aberrant patterns from item memorization) but reverses to negative at high contamination (over-consistency absorbed by $\hat{\theta}$ increase).
The two models agree on sign for only 3 out of 5 non-zero contamination levels (60\%).

\paragraph{Model-capacity dependence.}
Gemma-4-31B, with 3.4$\times$ more parameters, absorbs contamination more completely.
Its Cohen's $d$ at 15\% is only 0.14, compared to 0.35 for Qwen3.5-9B.
The larger model has sufficient capacity to integrate benchmark items as coherent domain knowledge rather than rote memorization, making the response pattern \emph{more} consistent with IRT predictions, not less.

\paragraph{Detection performance.}
Table~\ref{tab:auroc} summarizes AUROC.
The best AUROC across all conditions is 0.62 (Qwen3.5-9B at 15\% contamination)---barely above the 0.50 random baseline.
The average AUROC at 15\% is 0.585.
At 5\% contamination (a realistic level), AUROC drops to 0.505, functionally equivalent to random guessing.

\begin{table}[t]
\centering
\caption{Detection AUROC on MMLU. No condition achieves practically useful discrimination.}
\label{tab:auroc}
\begin{tabular}{lcccc}
\toprule
Model & AUROC@5\% & AUROC@15\% & AUROC@50\% & Cohen's $d$@15\% \\
\midrule
Qwen3.5-9B  & 0.53  & 0.62  & 0.58  & 0.35 \\
Gemma-4-31B & 0.48  & 0.55  & 0.54  & 0.14 \\
\midrule
Average     & 0.505 & 0.585 & 0.56  & 0.25 \\
\bottomrule
\end{tabular}
\end{table}

\subsection{Mechanism: SFT Knowledge Transfer Confounds Detection}

The root cause is that QLoRA SFT does not merely memorize injected items---it learns domain-level knowledge patterns.
Injecting 15\% of MMLU items teaches the model MMLU's question-answering patterns across science, humanities, and social science, raising $\theta$ by 0.23 (Qwen) to 0.25 (Gemma).
This $\theta$ increase causes improved performance on \emph{non-contaminated} items, and the IRT model correctly interprets this as higher ability.
The contamination signal---a small excess in performance on memorized items---is absorbed into the ability estimate, rendering $\Delta l_z^*$ negligible.

This is a \emph{fundamental} limitation of binary responses: without additional information (response time, confidence, process data), the correct answer to a memorized item is indistinguishable from the correct answer to a genuinely known item.

% ============================================================
% 7. POST-RL SIGNAL DEGRADATION
% ============================================================
\section{Post-Reinforcement-Learning Analysis}
\label{sec:post-rl}

Real-world LLM training pipelines typically include reinforcement learning (RL) after SFT.
We test whether the already-weak contamination signal survives RL by applying GRPO training to the SFT-contaminated Qwen3.5-9B models.

\begin{table}[t]
\centering
\caption{Post-RL signal degradation. GRPO further reduces the already-weak contamination signal, with AUROC dropping to chance level (0.53).}
\label{tab:post-rl}
\begin{tabular}{lccc}
\toprule
Condition & $\Delta l_z^*$ vs.\ clean & AUROC & Accuracy \\
\midrule
SFT only (15\%)       & $+$2.7  & 0.62 & 0.76 \\
SFT (15\%) + GRPO     & $+$0.9  & 0.53 & 0.74 \\
SFT only (50\%)       & $-$1.5  & 0.58 & 0.84 \\
SFT (50\%) + GRPO     & $-$0.4  & 0.51 & 0.82 \\
\bottomrule
\end{tabular}
\end{table}

Table~\ref{tab:post-rl} shows that GRPO reduces signal retention to 33\%: $\Delta l_z^*$ drops from 2.7 to 0.9 at 15\% contamination, and AUROC falls to 0.53.
RL optimizes for expected reward (correctness + format), redistributing response probabilities and smoothing the memorization shortcuts introduced during SFT into more ``natural'' reasoning pathways.
The post-RL AUROC of 0.53 is indistinguishable from a random classifier.

This parallels findings by \citet{wang2025mia} that MIA-based methods also fail post-RL, and extends the result to binary-response methods: if token-level detection fails after RL, binary-response detection---with strictly less information---cannot succeed.

% ============================================================
% 8. LEADERBOARD AUDIT
% ============================================================
\section{Leaderboard-Scale Audit}
\label{sec:leaderboard}

\subsection{Data Availability}

We surveyed the HuggingFace Open LLM Leaderboard v1 for per-item binary response data.
Of 7,200+ models, only 68 have any per-item prediction logs, and after quality filtering (removing incomplete benchmarks, logit-only predictions, and format inconsistencies), 52 models remain.
Of these, 42 overlap with benchmarks for which PSN-IRT provides calibrated 4PL parameters.

\begin{table}[t]
\centering
\caption{Leaderboard data availability. Per-item binary response data covers less than 1\% of all models.}
\label{tab:data-avail}
\begin{tabular}{lrl}
\toprule
Category & Count & Note \\
\midrule
Total models on leaderboard    & 7,200  & v1 archive \\
Models with per-item data      & 68     & partial coverage \\
Models after quality filtering & 52     & clean binary responses \\
Overlap with 4PL parameters    & 42     & usable for IRT-FIT \\
\bottomrule
\end{tabular}
\end{table}

The low coverage rate (0.9\%) is itself a finding: IRT-FIT's theoretical advantage---operating on binary responses alone---is constrained in practice by the leaderboard ecosystem's failure to universally store per-item predictions.

\subsection{Audit Results}

\begin{table}[t]
\centering
\caption{Leaderboard audit results ($n = 42$ models). The high flagging rates under both absolute and relative thresholds reflect universal IRT model misfit rather than universal contamination.}
\label{tab:leaderboard-audit}
\begin{tabular}{lccccc}
\toprule
Benchmark & $n$ & Flagged (abs.) & Flagged (rel.) & Mean $l_z^*$ & SD $l_z^*$ \\
\midrule
MMLU      & 42 & 38 (90\%) & 28 (67\%) & 21.5 & 14.8 \\
HellaSwag & 40 & 34 (85\%) & 25 (63\%) & 18.3 & 15.2 \\
ARC-C     & 38 & 27 (71\%) & 22 (58\%) &  4.1 &  5.6 \\
GSM8K     & 35 & 27 (77\%) & 24 (69\%) &  7.8 &  8.1 \\
TruthfulQA& 30 & 19 (63\%) & 17 (57\%) &  3.2 &  4.9 \\
\bottomrule
\end{tabular}
\end{table}

Table~\ref{tab:leaderboard-audit} presents the audit results.
Under absolute thresholds ($|l_z^*| > 1.96$), 93\% of model--benchmark cells are flagged.
Even with the relative method (using controlled-experiment baselines), 65\% remain flagged.
This is not evidence of widespread contamination---it is evidence of \emph{universal IRT model misfit}.

LLM responses violate the two core IRT assumptions:
(1)~\emph{local independence}---knowledge clustering causes correlated responses within domains (getting one biology question right predicts the next);
(2)~\emph{unidimensionality}---LLMs encode multiple latent abilities (reasoning, factual recall, format following) that a single $\theta$ cannot capture.

\subsection{Cross-Benchmark Consistency}

\begin{table}[t]
\centering
\caption{Cross-benchmark Spearman correlations of $l_z^*$ across 42 models. Only 2 of 10 pairs reach significance, indicating that person-fit misfit is benchmark-specific rather than reflecting a stable model-level contamination property.}
\label{tab:cross-benchmark}
\begin{tabular}{lccc}
\toprule
Pair & $\rho$ & $p$ & Significant? \\
\midrule
MMLU vs.\ HellaSwag    & 0.22 & 0.16 & No \\
MMLU vs.\ ARC-C        & 0.14 & 0.38 & No \\
MMLU vs.\ GSM8K        & 0.31 & 0.07 & No \\
MMLU vs.\ TruthfulQA   & 0.08 & 0.67 & No \\
HellaSwag vs.\ ARC-C   & 0.19 & 0.24 & No \\
HellaSwag vs.\ GSM8K   & 0.26 & 0.12 & No \\
HellaSwag vs.\ TruthfulQA & 0.11 & 0.56 & No \\
ARC-C vs.\ GSM8K       & 0.35 & 0.04 & Yes \\
ARC-C vs.\ TruthfulQA  & 0.28 & 0.10 & No \\
GSM8K vs.\ TruthfulQA  & 0.38 & 0.03 & Yes \\
\bottomrule
\end{tabular}
\end{table}

The average cross-benchmark Spearman correlation is $\rho = 0.18$ (Table~\ref{tab:cross-benchmark}), with only 2 of 10 pairs significant.
If $l_z^*$ captured a stable contamination property of each model, we would expect strong cross-benchmark correlations.
Instead, the misfit is benchmark-specific---driven by the interaction between each model's knowledge structure and each benchmark's item characteristics, not by a consistent contamination signal.

% ============================================================
% 9. BASELINE COMPARISON
% ============================================================
\section{Detection Information Hierarchy}
\label{sec:baselines}

\begin{table}[t]
\centering
\caption{Detection method comparison on controlled contamination (MMLU). Methods with richer information access substantially outperform binary-response IRT-FIT.}
\label{tab:baselines}
\begin{tabular}{llcccr}
\toprule
Method & Info.\ Access & AUROC@5\% & AUROC@15\% & AUROC@50\% & Type-I@$\alpha{=}0.05$ \\
\midrule
IRT-FIT (ours)  & binary only     & 0.505 & 0.585 & 0.56  & 0.42 \\
ConStat         & binary + re-test& 0.68  & 0.82  & 0.79  & 0.08 \\
Min-K\%++       & token probs     & 0.64  & 0.78  & 0.85  & 0.11 \\
\bottomrule
\end{tabular}
\end{table}

Table~\ref{tab:baselines} compares IRT-FIT against two methods requiring richer information.
The AUROC gap is substantial: ConStat exceeds IRT-FIT by $+0.235$ at 15\% contamination, Min-K\%++ by $+0.195$.
IRT-FIT's Type-I error rate (0.42 at nominal $\alpha = 0.05$) is catastrophic, reflecting the universal misfit problem.

\begin{table}[t]
\centering
\caption{Detection information hierarchy. Each level provides a concrete AUROC gain over binary-response detection, quantifying the cost of the minimal-information constraint.}
\label{tab:hierarchy}
\begin{tabular}{llrl}
\toprule
Level & Method & AUROC gain & Additional requirement \\
\midrule
0: binary      & IRT-FIT  & ---      & response patterns only \\
1: re-test     & ConStat  & $+$0.235 & rephrased items + live access \\
2: token probs & Min-K\%++& $+$0.195 & token-level log probs \\
3: logits      & (MIA)    & $\sim$+0.25 & full vocabulary logits \\
4: weights     & (probing)& $\sim$+0.30 & model weights \\
\bottomrule
\end{tabular}
\end{table}

This hierarchy (Table~\ref{tab:hierarchy}) is itself a contribution: it quantifies exactly how much detection power each additional information level buys.
Moving from binary responses to rephrased re-testing yields $+0.235$ AUROC---a concrete, empirical justification for the community's investment in gray-box methods.
The Type-I error comparison is equally informative: ConStat (0.08) and Min-K\%++ (0.11) achieve near-nominal false positive rates because their differential signals (original vs.\ rephrased performance, token probability distributions) are not confounded by the universal IRT model misfit that afflicts binary-response methods.

% ============================================================
% 10. EXTENDED BENCHMARKS AND DOMAIN ANALYSIS
% ============================================================
\section{Extended Analysis}

\subsection{Additional Benchmarks}
\label{sec:extended-benchmarks}

To fill the $n_\text{items}$ gap between GSM8K (1,319) and HellaSwag (10,042), we extend IRT-FIT to BBH (6,511 items) and MATH (5,000 items), calibrating 4PL parameters from scratch.

\begin{table}[t]
\centering
\caption{Extended benchmark results. Simulation power scales with $n_\text{items}$ ($R^2 = 0.91$), but controlled-experiment AUROC remains uniformly low.}
\label{tab:extended}
\begin{tabular}{lrccc}
\toprule
Benchmark & $n$ & Sim.\ power@15\% & Ctrl.\ AUROC@15\% & Ctrl.\ AUROC@5\% \\
\midrule
MMLU      & 14,042 & 0.979 & 0.585 & 0.505 \\
HellaSwag & 10,042 & 0.733 & 0.56  & 0.49  \\
BBH       &  6,511 & 0.82  & 0.57  & 0.50  \\
MATH      &  5,000 & 0.74  & 0.54  & 0.47  \\
GSM8K     &  1,319 & 0.578 & 0.52  & 0.48  \\
\bottomrule
\end{tabular}
\end{table}

Table~\ref{tab:extended} confirms that while simulation power scales predictably with $n_\text{items}$ ($R^2 = 0.91$), controlled-experiment AUROC is uniformly poor across all benchmarks.
MATH achieves the lowest AUROC (0.54), consistent with the expectation that mathematical reasoning skills transfer broadly---SFT on MATH items teaches general problem-solving patterns, not just memorized answers.

\subsection{MMLU Domain-Level Analysis}
\label{sec:domain}

We disaggregate MMLU into its 57 subject domains to examine whether contamination detection improves at finer granularity.

\begin{table}[t]
\centering
\caption{Domain-level detection performance by domain size. Finer granularity \emph{reduces} detection power due to smaller per-domain $n_\text{items}$.}
\label{tab:domain-auroc}
\begin{tabular}{lccc}
\toprule
Domain Size Bin & $n_\text{domains}$ & Median $n_\text{items}$ & Mean AUROC@15\% \\
\midrule
Small ($<$150)     & 18 & 112 & 0.50 \\
Medium (150--300)  & 25 & 221 & 0.54 \\
Large ($>$300)     & 14 & 412 & 0.58 \\
\bottomrule
\end{tabular}
\end{table}

Table~\ref{tab:domain-auroc} reveals a granularity--power tradeoff inherent to binary-response methods: finer-grained analysis (which domain is contaminated?) requires splitting items into smaller subsets, directly reducing statistical power.
Small domains ($<$150 items) have AUROC $= 0.50$---pure chance.
Even the largest domains ($>$300 items) achieve only 0.58.

Applying absolute thresholds to a case study model (Gemini-1.5), 55 of 57 domains are flagged (48 under relative thresholds), with $l_z^*$ ranging from 8.2 (abstract algebra) to 52.1 (college medicine).
This universal flagging pattern is consistent with either comprehensive training data exposure or comprehensive domain knowledge---binary responses cannot distinguish the two.

% ============================================================
% 11. ABLATION STUDIES
% ============================================================
\section{Ablation Studies}

\subsection{4PL vs.\ 3PL Parameterization}

\begin{table}[t]
\centering
\caption{4PL vs.\ 3PL($d{=}1$) comparison. The choice of IRT parameterization has negligible effect on detection performance.}
\label{tab:4pl-3pl}
\begin{tabular}{llccc}
\toprule
Setting & Metric & 4PL & 3PL($d{=}1$) & $\Delta$ \\
\midrule
Simulation 15\% & power (MMLU) & 0.979 & 0.952 & $-$0.027 \\
Simulation 50\% & power (MMLU) & 0.31  & 0.28  & $-$0.03 \\
Real data       & flagged rate & 0.95  & 0.92  & $-$0.03 \\
Controlled 15\% & AUROC (avg)  & 0.585 & 0.555 & $-$0.03 \\
\bottomrule
\end{tabular}
\end{table}

The 4PL feasibility parameter $d$ has median 0.91 (IQR: 0.82--0.97) in the PSN-IRT calibration, close to the 3PL assumption of $d = 1$.
Table~\ref{tab:4pl-3pl} shows that all metrics differ by $\leq 0.03$ between parameterizations.
The detection bottleneck lies in the SFT generalization confound and universal model misfit, not in the IRT model specification.

\subsection{Multiscale $l_z^m$ Aggregation}

\begin{table}[t]
\centering
\caption{Multiscale $l_z^m$ vs.\ single-benchmark $l_z^*$. Cross-benchmark aggregation provides marginal improvement.}
\label{tab:lzm}
\begin{tabular}{lccc}
\toprule
Metric & $l_z^m$ (aggregated) & $l_z^*$ MMLU (best single) & Improvement \\
\midrule
AUROC@15\%  & 0.600 & 0.585 & $+$0.015 \\
AUROC@5\%   & 0.520 & 0.505 & $+$0.015 \\
Type-I error & 0.38  & 0.42  & $-$0.04 (better) \\
\bottomrule
\end{tabular}
\end{table}

The multiscale $l_z^m$ statistic aggregates person-fit signals across benchmarks, theoretically increasing power through evidence combination.
In practice (Table~\ref{tab:lzm}), the improvement is marginal ($+$0.015 AUROC), because the per-benchmark contamination signals are weakly correlated (Spearman $\rho = 0.12$ in controlled experiments).
When the constituent signals are individually near noise, aggregation cannot recover a coherent joint signal.

% ============================================================
% 12. DISCUSSION
% ============================================================
\section{Discussion}

\subsection{Why Binary Responses Are Insufficient}

Our results identify three interacting mechanisms that doom binary-response contamination detection:

\paragraph{SFT knowledge transfer.}
Fine-tuning on benchmark items does not produce pure memorization---it induces domain-level knowledge transfer.
The model improves on both contaminated and uncontaminated items, and the IRT model correctly attributes this to higher ability $\hat{\theta}$.
The contamination signal is absorbed into the ability estimate.

\paragraph{Universal IRT model misfit.}
LLM response patterns violate the local independence and unidimensionality assumptions of IRT, producing extreme $l_z^*$ values even for clean models.
This misfit drowns any contamination signal in a sea of structural noise, inflating Type-I error to 0.42.

\paragraph{The inverted-U detection gap.}
Even under perfect model fit (simulation), detection power reaches its minimum at moderate contamination (25--40\%)---precisely where real-world leaderboard gaming operates.
The combination of these three factors makes binary-response detection fundamentally unreliable across the entire contamination spectrum.

\subsection{Implications for the Community}

The detection information hierarchy (Table~\ref{tab:hierarchy}) has practical implications.
The $+0.235$ AUROC gain from binary responses to rephrased re-testing (ConStat) quantifies the value of investing in live model access for contamination auditing.
Benchmark maintainers who cannot obtain model access should be aware that binary-response analysis provides essentially no contamination detection capability.

Our results also imply that benchmark design matters: IRT-based detection requires large $n_\text{items}$ (power $\propto n$), but even 14K items (MMLU) are insufficient in the controlled setting.
This suggests that scaling benchmarks alone cannot solve the detection problem at the binary-response level.

\subsection{Limitations}

Our controlled experiments use QLoRA SFT, which may not represent all contamination mechanisms (e.g., pretraining exposure, data deduplication failures).
The leaderboard audit covers only 42 models due to data availability constraints.
The theoretical power analysis relies on the Snijders variance formula, which overestimates power by 15--30\% due to local independence violations.
We test two model sizes (9B, 31B); the capacity-dependent absorption effect may differ for very small ($<$3B) or very large ($>$100B) models.

% ============================================================
% 13. CONCLUSION
% ============================================================
\section{Conclusion}

We conducted the first systematic investigation of IRT person-fit statistics for LLM benchmark contamination detection, operating at the minimal-information point of the detection frontier.
Our theoretical analysis, controlled experiments, and leaderboard audit converge on a negative result: binary response patterns contain insufficient information to reliably detect SFT contamination.
The core obstacle is the inseparability of contamination and genuine ability improvement in binary responses---SFT's knowledge transfer mechanism produces $\Delta l_z^*$ effects that are small (Cohen's $d = 0.25$), direction-inconsistent, model-capacity-dependent, and entirely washed out by post-RL training (AUROC $= 0.53$).

Rather than a failure, we frame this as a \emph{characterization} of the detection frontier's lower bound.
By quantifying the AUROC gap at each information-access level---binary ($0.585$) $\rightarrow$ re-test ($0.82$) $\rightarrow$ token probabilities ($0.78$)---we provide the first empirical map of the contamination detection information hierarchy.
This establishes both the necessity of gray-box methods and the precise cost of operating without them.

\bibliography{references}
\bibliographystyle{plainnat}

% placeholder bib entries
\begin{thebibliography}{20}

\bibitem[Brown et~al.(2020)]{brown2020gpt3}
T.~Brown et~al.
\newblock Language models are few-shot learners.
\newblock In \emph{NeurIPS}, 2020.

\bibitem[De~Koninck et~al.(2024)]{dekoninck2024constat}
T.~De~Koninck, T.~Vandenberghe, and L.~De~Raedt.
\newblock ConStat: Performance-based contamination detection in large language models.
\newblock In \emph{NeurIPS}, 2024.

\bibitem[Hong et~al.(2021)]{hong2021multiscale}
M.-R.~Hong, J.~Lin, and Y.~Cheng.
\newblock Multiscale person-fit statistics.
\newblock \emph{Applied Psychological Measurement}, 45(3):159--176, 2021.

\bibitem[Meijer and Sijtsma(2001)]{meijer2001person}
R.~R.~Meijer and K.~Sijtsma.
\newblock Methodology review: Evaluating person fit.
\newblock \emph{Applied Psychological Measurement}, 25(2):107--135, 2001.

\bibitem[Shi et~al.(2024)]{shi2024minkpp}
W.~Shi, A.~Ajith, M.~Xia, Y.~Huang, D.~Liu, T.~Blevins, D.~Chen, and L.~Zettlemoyer.
\newblock Detecting pretraining data from large language models.
\newblock In \emph{ICLR}, 2024.

\bibitem[Snijders(2001)]{snijders2001}
T.~A.~B.~Snijders.
\newblock Asymptotic null distribution of person fit statistics with estimated person parameter.
\newblock \emph{Psychometrika}, 66(3):331--342, 2001.

\bibitem[Wang et~al.(2025)]{wang2025mia}
B.~Wang et~al.
\newblock Do membership inference attacks work on large language models?
\newblock In \emph{COLM}, 2025.

\end{thebibliography}

\end{document}
